% Section 5 — DocRED Results (Draft)
% Data source: Phase 2 DocRED registry experiments (checker_version: fixed_20260514)
% exp_ids: phase2_docred_unconstrained, phase2_docred_grammar_only,
%          phase2_docred_semantic_only, phase2_docred_full_s1
% All numbers verified against registry 2026-05-14.
% F1 computed via all_set method with name→Wikidata PID normalization.

\subsection{Cross-Dataset Validation on DocRED}
\label{sec:docred-results}

Table~\ref{tab:docred-ablation} presents the four-way ablation on DocRED,
a document-level relation extraction benchmark with 96 Wikidata relation
types---14$\times$ the ontology size of SciERC
(Section~\ref{sec:scierc-results}).
All experiments use Qwen3-8B in a zero-shot setting on 988 test documents
with identical prompts and decoding parameters.
Relation F1 is computed via name-to-Wikidata-PID normalization:
each predicted relation name is mapped to its canonical PID before
matching against gold annotations, and predictions that cannot be mapped
to any valid PID are counted as false positives.

% ──────────────────────────────────────────
%  Table: DocRED four-way ablation
% ──────────────────────────────────────────
\begin{table}[t]
\centering
\small
\setlength{\tabcolsep}{3.6pt}
\begin{tabular}{@{}l rrrrr rr@{}}
\toprule
& \multicolumn{5}{c}{\textbf{Violation Rates (\%)}}
& \multicolumn{2}{c}{\textbf{F1}} \\
\cmidrule(lr){2-6} \cmidrule(l){7-8}
\textbf{Mode}
  & DR\,$\downarrow$ & $\Delta$DR
  & RV\,$\downarrow$ & ET\,$\downarrow$ & CF
  & Rel.\,$\uparrow$ & Ent.\,$\uparrow$ \\
\midrule
Unconstrained
  & 18.46 & ---
  & 29.66 & 8.18
  & 0.10
  & 7.74  & 58.15 \\
Grammar
  & 26.17 & $\uparrow$42\%
  & \textbf{0.03} & \textbf{0}
  & 0.07
  & \textbf{8.31}  & --- \\
Semantic
  & 13.84 & $\downarrow$25\%
  & 29.87 & 10.44
  & 1.05
  & 7.70  & 57.81 \\
Full (S1)
  & \textbf{10.19} & $\downarrow$\textbf{45\%}
  & \textbf{0.03} & \textbf{0}
  & 2.57
  & 8.23  & \textbf{59.37} \\
\bottomrule
\end{tabular}
\caption{Four-way ablation on DocRED (988 documents, Qwen3-8B, zero-shot).
DR: domain/range violation rate;
RV: relation vocabulary violation rate;
ET: entity type violation rate;
CF: cross-field consistency violation rate.
$\Delta$DR: relative change vs.\ the unconstrained baseline
($\uparrow$: increase; $\downarrow$: decrease).
Best per column in \textbf{bold} (CF excluded---see text).
F1 scores are multiplied by 100.
Grammar-only entity F1 is unavailable (---).}
\label{tab:docred-ablation}
\end{table}

% ──────────────────────────────────────────
%  (a) Core result
% ──────────────────────────────────────────
\paragraph{Core result: DR halved at zero F1 cost.}
The full two-layer configuration reduces domain/range violations from
18.46\% to 10.19\%, a 44.8\% relative reduction, while simultaneously
eliminating entity type violations (8.18\%\,$\to$\,0\%) and
near-eliminating relation vocabulary violations
(29.66\%\,$\to$\,0.03\%).
% phase2_docred_unconstrained: DR=0.1846, ET=0.0818, RV=0.2966
% phase2_docred_full_s1: DR=0.1019, ET=0, RV=0.0003
Relation F1 moves from 7.74 to 8.23 (+0.49\,pp) and entity F1 from
58.15 to 59.37 (+1.22\,pp);
% phase2_docred_unconstrained: R-F1=0.0774, E-F1=0.5815
% phase2_docred_full_s1: R-F1=0.0823, E-F1=0.5937
as with SciERC, these differences are within noise for single-seed
evaluation and are best characterized as F1-neutral.
Cross-field consistency violations increase modestly from 0.10\% to
2.57\%, substantially lower than the 9.90\% observed on SciERC;
% phase2_docred_unconstrained: CF=0.0010; phase2_docred_full_s1: CF=0.0257
we discuss this trade-off in Section~\ref{sec:discussion}.

% ──────────────────────────────────────────
%  (b) Grammar layer
% ──────────────────────────────────────────
\paragraph{Grammar layer: structural guarantees generalize, but surface latent violations.}
Grammar-only constraints generalize to DocRED's 96-relation ontology,
eliminating entity type violations (8.18\%\,$\to$\,0\%) and
near-eliminating relation vocabulary violations
(29.66\%\,$\to$\,0.03\%).
% phase2_docred_grammar_only: ET=0, RV=0.0003
However, domain/range violations \emph{increase} from 18.46\% to
26.17\% (+42\% relative).
% phase2_docred_grammar_only: DR=0.2617
This counter-intuitive result has a structural explanation.
In the unconstrained baseline, 29.66\% of triples contain
out-of-vocabulary relation names and are classified as RV violations;
because no schema entry exists for an invalid relation name,
these triples \emph{cannot} be evaluated against domain/range rules
and are therefore absent from the DR count.
When grammar constraints force all relation names into the valid
vocabulary, these formerly RV-classified triples receive valid
relation names---but the associated entity types frequently violate
the domain/range rules of the assigned relation, surfacing as new
DR violations.
The grammar layer thus \emph{reveals} latent type incompatibilities
rather than creating them, directly motivating the semantic layer:
without domain/range enforcement, grammar alone makes the type
violation landscape worse, not better.

% ──────────────────────────────────────────
%  (c) Enabling dependency
% ──────────────────────────────────────────
\paragraph{Enabling dependency confirmed across datasets.}
Semantic constraints alone reduce DR by 25\%
(18.46\%\,$\to$\,13.84\%) but cannot eliminate relation vocabulary
violations (RV\,=\,29.87\%) or entity type violations
(ET\,=\,10.44\%).
% phase2_docred_semantic_only: DR=0.1384, RV=0.2987, ET=0.1044
The elevated ET under semantic-only decoding (10.44\% vs.\ 8.18\%
unconstrained) arises because the domain/range mask blocks entity type
tokens at \texttt{head\_type}/\texttt{tail\_type} positions without a
grammar layer to constrain output to the type enum.
On DocRED, where entity types are frequently excluded from relations
(e.g., \textsc{NUM} is blocked in 92\% of relations), the model
escapes to OOV strings (e.g., ``ORGB'', ``EVENT''), producing 355
additional ET violations exclusively in relation type fields while
entity-list violations remain unchanged.
On SciERC, every entity type is a valid head and tail for all
relations, so the mask is never triggered and ET stays at 0.40\%.
Combining both layers (Full S1) reduces DR a further 26\% relative to
the semantic-only configuration (13.84\%\,$\to$\,10.19\%),
% phase2_docred_full_s1: DR=0.1019; delta=(13.84-10.19)/13.84=0.264
replicating the enabling dependency observed on SciERC:
without grammar-enforced vocabulary constraints, the semantic layer
encounters out-of-vocabulary relation types for which no domain/range
rules exist, limiting its enforcement capacity.
The consistency of this pattern across two datasets with substantially
different ontology sizes (7 vs.\ 96 relation types) and violation
profiles (SciERC RV\,=\,1.39\% vs.\ DocRED RV\,=\,29.66\%) confirms
that the grammar$\to$semantic enabling dependency is a general
architectural property of the framework, not an artifact of a
particular schema.

% ──────────────────────────────────────────
%  (d) Constraint tightness scaling
% ──────────────────────────────────────────
\paragraph{Constraint tightness governs enforcement ceiling.}
The residual DR on DocRED (10.19\%) substantially exceeds that of
SciERC (0.93\%, Section~\ref{sec:scierc-results}).
This gap is explained by per-relation \emph{constraint tightness}:
DocRED relations admit on average only 5.68 valid entity type pairs
(mean tightness\,=\,0.842, i.e., 84.2\% of entity type pairs violate
domain/range rules), compared to 23 valid pairs on SciERC
(tightness\,=\,0.333).
% DocRED: avg_valid_pairs=5.68, avg_tightness=0.842
% SciERC: avg_valid_pairs=23, avg_tightness=0.333
With tighter per-relation constraints, the semantic layer's logit mask
eliminates a larger fraction of candidate tokens at each decoding step,
increasing the probability that the intersection of grammar and
semantic masks yields a narrow or empty candidate set.
In the limiting case, the system falls back to grammar-only constraints,
preserving structural validity but forgoing domain/range enforcement
for that token position.
Conversely, SciERC's looser constraints leave more valid alternatives
for the model to redirect probability mass toward, enabling
near-complete DR elimination.
The DR reduction thus scales inversely with constraint tightness:
45\% on DocRED (tight) vs.\ 94\% on SciERC (loose).
% DocRED: DR_reduction=44.8%; SciERC: DR_reduction=93.9%
This represents, to our knowledge, the first empirical characterization
of how ontology constraint properties modulate decoding-time enforcement
effectiveness; we develop this analysis further in
Section~\ref{sec:discussion}.
